质量分别 \(m_1,m_2>0\) 的两光滑硬球在水平无外力空间发生一次短时碰撞。碰撞瞬间从球2指向球1的单位法向量为 \(\mathbf n\)，碰前速度为 \(\mathbf v_1,\mathbf v_2\)，相对法向速度 \(g_n=(\mathbf v_1-\mathbf v_2)\cdot\mathbf n<0\)；光滑无切向冲量，法向恢复系数 \(e\in[0,1]\)。记约化质量 \(\mu=m_1m_2/(m_1+m_2)\)。
(1) 利用动量守恒与恢复系数定义求冲量及两球碰后速度矢量。
(2) 证明碰后相对切向速度不变而法向速度反号乘 \(e\)，计算动能损失。
(3) 写出质心速度与质心系碰撞前后相对动能的表达，讨论 \(e=1\) 与 \(e=0\)。
(4) 对实验室中2球初始静止、1球以速率 \(v\) 入射的特例，求法向冲量、1球速度的两个分量，并检查一维正碰极限。
